\scanpage{163}
\mspage{112}
% archive sheet number: 163
et \unsure{incluant le|en dehors du} cas des imm. ouvertes et des imm. fermées ?

\uline{Démonstration du} th.~3 \deleted{\unsure{partie I}}\inserted{(p.~107)} Il suffit de prouver que si $i$ est \deleted{\unsure{$d$}}$W$-dewyérienne (\unsure{i.e.} $\unsure{i_0|i^*_0|i^*}$ \unsure{strictement} $W$-asphérique), alors $i'$ l'est également -- car il en résultera que $i'$ est $W$-asphérique et a fortiori $\in W$.

Il faut donc vérifier que pour \deleted{\unsure{$t$}} $t'\in\dot U'=\dot T$, l'élément $i'{}^*(t')$ \inserted{$=(u'\mapsto\operatorname{Hom}_{\unsure{X'|X}}(u',t'))$} de $U'{}^\wedge$ est \unsure{totalement} $W$-asphérique. Or $t'$ provient d'un unique élément $t$ de $\dot T$, et on sait qu'on a

\[
i'{}^*(\unsure{t|t'})=p_!(i^*(t)),\quad\text{où }p_!:U^\wedge\unsure{-|\to}U'{}^\wedge .
\]

Par hypothèse $F=i^*(t)$ est un élément tot. $W$-asphérique de $U^\wedge$, ce qui signifie que pour toute flèche $p:U-U'$ \inserted{dans $\mathrm{Cat}$} de source $U$, $p_!(F)$ est un élément $W$-asphérique de $\unsure{U'|U'{}^\wedge}$. En fait, on a le lemme suivant (par transitivité $(p'p)_!\simeq p'_!p_!$)

\uline{Lemme} \unsure{Soit}\deleted{\unsure{$F$}} $U$ dans $\mathrm{Cat}$, et \unsure{$F$|$R$} de $U^\wedge$. Si $F$ est tot. $W$-asph. \unsure{(resp.\textsuperscript{$W$}} ($U_{/F}$ est tot\textsuperscript{t} $W$-asphérique), alors pour tte flèche $p:U-U'$ dans $\mathrm{Cat}$ de source $U$, $F'=p_!(F)$ est un \unsure{objet} tot\textsuperscript{t} $W$-asph. de $U'{}^\wedge$ (resp.
